% SPDX-License-Identifier: CC-BY-4.0
% Copyright (c) 2026 Martin Schröder <info@swedishembedded.com>

\section{Related work}

SWE-bench poses resolving real GitHub issues as the task, with 2{,}294 instances graded
by the repository's tests \citep{jimenez2024swebench}. SWE-agent shows that the
interface between the model and the repository (file viewer, editor, shell) changes
what a model achieves \citep{yang2024sweagent}; we take the corresponding lesson that the
interaction, not only the final patch, is what must be recorded and trained on.

Three recent works build \emph{training} environments. SWE-Gym provides 2{,}438 Python
instances with executable environments and reports resolve-rate gains on SWE-bench
Verified and Lite from fine-tuned agents \citep{pan2025swegym}. R2E-Gym generates more than
8{,}700 tasks procedurally, using test generation and back-translation from commits
\citep{jain2025r2egym}. SWE-smith synthesizes task instances by modifying code in
a repository so that tests fail, building 50k instances from 128 repositories and
training a model it reports at 40.2\% pass@1 on SWE-bench Verified \citep{yang2025swesmith}.
All three, as described in their abstracts, pool repositories and evaluate on a fixed
external benchmark.

\paragraph{What differs here.} Our target is one repository's own domain, with the
adapter evaluated on that repository's families reserved before rollouts, and then
the \emph{sequence}: further repositories and later revisions, with retention measured on
everything met before. We borrow SWE-smith's mutation-and-validate idea and R2E-Gym's
commit-based idea as task producers; the learning loop, evidence store, release gate and
agent are the existing ones of Splinter, Sven and Brain, and the trainer is Brain's native
LoRA \citep{hu2021lora}.

\paragraph{Limits of this review.} We read the abstracts of six works and verified
their bibliographic data against arXiv; we did not read their evaluation sections in
full, so we make no claim that no prior work studies repository sequences or retention.
The claim of this paper is the measured answer to our question for our system, not
priority.
